% Generated draft: requires mathematical and source review.
\documentclass[11pt]{article}
\usepackage[margin=25mm]{geometry}
\usepackage{fontspec}
\setmainfont{Latin Modern Roman}
\usepackage{amsmath,amssymb,mathtools,unicode-math}
\setmathfont{Latin Modern Math}
\usepackage{enumitem,needspace,xurl,hyperref}
\setlength{\parindent}{0pt}
\setlength{\parskip}{5pt}
\setlength{\emergencystretch}{3em}
\newcommand{\N}{\mathbb{N}}
\newcommand{\Z}{\mathbb{Z}}
\newcommand{\Q}{\mathbb{Q}}
\newcommand{\R}{\mathbb{R}}
\newcommand{\C}{\mathbb{C}}
\newcommand{\F}{\mathbb{F}}
\newcommand{\A}{\mathbb{A}}
\newcommand{\PP}{\mathbb{P}}
\newcommand{\E}{\mathbb{E}}
\newcommand{\eps}{\varepsilon}
\newcommand{\dd}{\,\mathrm{d}}
\begin{document}
\section*{UnsolvedMath ID: 30002051}
\textbf{Source problem code:} catalogue entry
\subsection{Definitions and mathematical statement}
For a real number \(t\), write
\[
\lVert t\rVert=\min_{m\in\Z}|t-m|
\]
for its distance to the nearest integer. Interpreting the supplied term ``pseudo-absolute-value sequence'' as a sequence \(D=(n_k)_{k\geq0}\) of positive integers, suppose that
\[
n_0=1,\qquad n_k\mid n_{k+1}\quad(k\geq0).
\]
For each positive integer \(q\), define
\[
|q|_D=\min\bigl\{n_k^{-1}:k\geq0,\ q\in n_k\Z\bigr\}.
\]
Under this interpretation the minimum exists: the set is nonempty because \(n_0=1\), and its distinct values arise from positive divisors of \(q\).

The assertion supplied under the title ``Mixed Littlewood Conjecture for Pseudo-Absolute Values'' is
\[
\text{for every such }D\text{ and every }x\in\R,
\qquad
\inf_{q\in\Z_{\geq1}} |q|_D\,\lVert qx\rVert=0.
\]
\textbf{Unresolved qualifications.} The record does not explicitly specify the domain of the sequence terms or require strict growth or unboundedness. Positive integer terms are an explicit interpretive assumption in this draft; no growth condition has been added. The displayed expression is retained exactly as supplied and contains no multiplicative factor \(q\). Its intended normalization and its identification with the named conjecture require checking against the report. The supplied verification label does not resolve these issues independently.
\subsection{Short English statement}
For each admissible divisibility sequence, the supplied question asks whether the product of its reciprocal-divisor weight and the distance of \(qx\) to the nearest integer has infimum zero over positive integers \(q\), for every real \(x\). The sequence convention and the normalization of the displayed formula require human source review.
\subsection{Sources}
\begin{itemize}
\item Supplied local UnsolvedMath registry record, permanent ID \(30002051\), source code \texttt{catalogue entry}, in the Oberwolfach Reports Open Problems collection; registry reference: \texttt{data/unsolvedmath/problems.json}. No individual proposer is supplied.
\item Supplied report citation: \emph{Diophantische Approximationen (2013)}, Oberwolfach Reports, pp.~1307--1359, DOI \texttt{10.4171/owr/2012/22}: \url{https://doi.org/10.4171/owr/2012/22}. The record describes these reports as documenting mathematical workshops and open-problem discussions. Its proposed year and year tag are 2012, whereas its citation title gives 2013; this discrepancy is unresolved.
\item Harrap--Haynes, \emph{The mixed Littlewood conjecture for pseudo-absolute values}: \url{https://arxiv.org/abs/1012.0191}.
\item Liu, \emph{Some refined results on the mixed Littlewood conjecture}: \url{https://doi.org/10.1017/S1446788718000198}.
\item Supplied dated literature curation, checked \texttt{2026-08-21}: status ``open.'' The supplied assessment says that mixed Littlewood statements for special \(p\)-adic or pseudo-absolute sequences and for almost every parameter are known, while the universal assertion for every pseudo-absolute-value sequence remains conjectural. This is supplied dated curation derived from the OWR material, not independently verified; its time-sensitive status, cited claims, and applicability to the displayed normalization require review.
\item The supplied statement status is \texttt{corrected\_verified}, with the assertion that the official report states the conjecture in this form. This is a source-record claim, not an independent verification performed for this draft.
\end{itemize}
\end{document}
